% =====================================================================
%  requerimientos.tex  --  Documento de Requerimientos de LEIP
%  ARCHIVO UNICO: no necesita ningun otro archivo.
%  Compilar en Overleaf con pdfLaTeX (dos pasadas por el indice).
% =====================================================================
\documentclass[11pt,a4paper]{article}

% ---------------------------------------------------------------------
%  DATOS DEL PROYECTO  (tocá solo esto para actualizar el documento)
% ---------------------------------------------------------------------
\newcommand{\proyectoNombre}{LEIP}
\newcommand{\proyectoNombreLargo}{LatAm Economic Intelligence Platform}
\newcommand{\proyectoUniversidad}{Universidad Cat\'olica del Uruguay}
\newcommand{\proyectoFacultad}{Facultad de Ingenier\'ia y Tecnolog\'ias}
\newcommand{\proyectoCarrera}{Ingenier\'ia en Inform\'atica}
\newcommand{\docVersion}{1.1}
\newcommand{\docTitulo}{Especificaci\'on de Requerimientos}
\newcommand{\docFecha}{Septiembre 2026}
\newcommand{\modMacro}{Macro Expectations Intelligence}
\newcommand{\modPrecios}{Prices \& Inflation Intelligence}
\newcommand{\modActividad}{Real Activity Intelligence}

% ---------------------------------------------------------------------
%  PREAMBULO  (estilos y comandos; normalmente no hace falta tocarlo)
% ---------------------------------------------------------------------
\usepackage[utf8]{inputenc}
\usepackage[T1]{fontenc}
\IfFileExists{lmodern.sty}{\usepackage{lmodern}\usepackage{microtype}}{}

\IfFileExists{spanish.ldf}{%
  \usepackage[spanish,es-tabla,es-noquoting]{babel}%
}{%
  \usepackage[english]{babel}%
  \renewcommand{\contentsname}{\'Indice}%
  \renewcommand{\listfigurename}{\'Indice de figuras}%
  \renewcommand{\listtablename}{\'Indice de tablas}%
  \renewcommand{\tablename}{Tabla}%
  \renewcommand{\figurename}{Figura}%
  \renewcommand{\refname}{Referencias}%
}

\usepackage[a4paper,margin=2.5cm,headheight=15pt]{geometry}
\usepackage{fancyhdr}
\usepackage{parskip}
\usepackage{booktabs}
\usepackage{tabularx}
\usepackage{longtable}
\usepackage{array}
\usepackage{enumitem}
\setlist{noitemsep,topsep=4pt}
\usepackage{graphicx}
\usepackage{xcolor}
\usepackage{tcolorbox}
\tcbuselibrary{breakable}
\usepackage[hidelinks]{hyperref}
\usepackage{url}
\urlstyle{same}

\definecolor{leipAzul}{HTML}{1F4E79}
\definecolor{leipGris}{HTML}{5A5A5A}
\definecolor{leipAmbar}{HTML}{B8860B}
\definecolor{leipAmbarBg}{HTML}{FFF8E1}

\pagestyle{fancy}
\fancyhf{}
\fancyhead[L]{\small\color{leipGris}\proyectoNombre{} --- \docTitulo}
\fancyhead[R]{\small\color{leipGris} v\docVersion}
\fancyfoot[C]{\small\thepage}
\renewcommand{\headrulewidth}{0.4pt}

\usepackage{titlesec}
\titleformat{\section}{\Large\bfseries\color{leipAzul}}{\thesection}{0.6em}{}
\titleformat{\subsection}{\large\bfseries\color{leipAzul}}{\thesubsection}{0.6em}{}
\titleformat{\subsubsection}{\normalsize\bfseries\color{leipGris}}{\thesubsubsection}{0.6em}{}

% --- Portada -----------------------------------------------------------
\newcommand{\logoUniversidad}{img/logo-ucu.png}
\newcommand{\portada}[1]{%
  \begin{titlepage}
  \centering
  \vspace*{1.5cm}
  \IfFileExists{\logoUniversidad}{%
    \includegraphics[width=0.40\textwidth]{\logoUniversidad}\par
    \vspace{1.2cm}%
  }{%
    \fbox{\parbox[c][3cm][c]{0.5\textwidth}{\centering\small\texttt{\logoUniversidad}}}\par
    \vspace{1.2cm}%
  }
  {\Large \proyectoUniversidad\par}
  {\large \proyectoFacultad\par}
  {\large \proyectoCarrera\par}
  \vspace{3cm}
  {\Huge\bfseries\color{leipAzul} #1\par}
  \vspace{0.8cm}
  {\Large \proyectoNombre{} --- \proyectoNombreLargo\par}
  \vspace{2cm}
  {\large Versi\'on \docVersion\par}
  {\large \docFecha\par}
  \vfill
  \begin{tabular}{c}
    Sebastian Forische \\
    Armando Hernandez de la Vega \\
    Luana Acu\~na \\
    Bruno Sebastian Olivera Viera Gomez \\
  \end{tabular}
  \vspace{1.5cm}
  \end{titlepage}
}

% --- Historial de revisiones -------------------------------------------
\newenvironment{historial}{%
  \section*{Historial de revisiones}
  \addcontentsline{toc}{section}{Historial de revisiones}
  \begin{tabular}{@{}l l p{0.52\textwidth} l@{}}
  \toprule
  \textbf{Versi\'on} & \textbf{Fecha} & \textbf{Descripci\'on} & \textbf{Autor} \\
  \midrule
}{%
  \bottomrule
  \end{tabular}
  \clearpage
}

% --- Requerimientos funcionales y no funcionales -----------------------
\newcommand{\rf}[2]{%
  \noindent
  \textbf{RF-#1:} #2
  \par\smallskip
}

\newcommand{\rnf}[2]{%
  \noindent
  \textbf{RNF-#1:} #2
  \par\smallskip
}

% =====================================================================
%  DOCUMENTO
% =====================================================================

\begin{document}

\portada{\docTitulo}

\begin{historial}
1.0 & 31/08/2026 & Versión inicial. Extracción de requerimientos desde el Documento de Visión & Armando \\
1.1 & 11/09/2026 & Se agregan introducción, alcance y convenciones. Se retira el término MVP y se reformulan los requerimientos que dependían de nombres comerciales de planes & Equipo \\
\end{historial}

\tableofcontents
\clearpage

% =====================================================================
\section{Introducción}

\subsection{Propósito}

Este documento especifica, de forma numerada y verificable, qué debe hacer
\proyectoNombre{} y bajo qué condiciones debe hacerlo. Traduce a requerimientos
concretos las capacidades y restricciones descritas en el Documento de Visión, y
constituye la referencia contra la cual se especifican los casos de uso, se
diseña la arquitectura y se verifica el producto.

Este documento es la versión vigente de los requerimientos. No se considera
cerrado: se actualiza cada vez que se agrega, modifica o retira un
requerimiento, y cada cambio queda asentado en el historial de revisiones.

\subsection{Alcance del documento}

El documento cubre los requerimientos funcionales del sistema de datos, la capa
de inteligencia y la aplicación web, y los requerimientos no funcionales de
trazabilidad, seguridad, aislamiento y cumplimiento.

No cubre la descripción paso a paso de las interacciones, que corresponde a la
Especificación de Casos de Uso, ni las decisiones técnicas con las que se
satisfacen estos requerimientos, que corresponden a las Notas de Arquitectura y
Diseño. Tampoco cubre los aspectos comerciales del producto, que se desarrollan
en el Documento de Negocio.

\subsection{Convención de identificadores}

Los requerimientos se identifican como \textbf{RF-nn} (funcionales) y
\textbf{RNF-nn} (no funcionales). Estos identificadores son estables: una vez
asignados no se reutilizan ni se renumeran, aunque el requerimiento se elimine.
Los casos de uso, casos de prueba y decisiones de arquitectura referencian estos
identificadores.

% =====================================================================
\section{Requerimientos funcionales}
\label{sec:rf}

Los requerimientos funcionales describen qué debe hacer el sistema. Se agrupan
en tres bloques que siguen el recorrido de la información dentro del producto:
el sistema de datos, que obtiene y normaliza la información oficial; la capa de
inteligencia, que calcula sobre ella; y la aplicación web, a través de la cual
el usuario accede al resultado.

% ---------------------------------------------------------------------
\subsection{Sistema de datos}

Este bloque cubre el recorrido de un dato desde la publicación del organismo
emisor hasta quedar disponible para el cálculo: cómo se obtiene, cómo se lleva
al modelo común y qué controles atraviesa antes de publicarse.

\subsubsection{Ingesta de fuentes oficiales}

\rf{ingesta-obtener}{El sistema debe obtener informaci\'on econ\'omica oficial
para los m\'odulos incluidos en la primera etapa, desde las fuentes definidas
para cada pa\'is.}

\rf{ingesta-programada}{El sistema debe ejecutar actualizaciones autom\'aticas
seg\'un la frecuencia de publicaci\'on de cada fuente.}

\rf{ingesta-registro}{El sistema debe registrar, para cada dato o archivo
incorporado: fuente, fecha de publicaci\'on, fecha de ingesta, formato,
versi\'on y \emph{checksum}.}

\subsubsection{Normalización de datos}

\rf{norm-modelo}{El sistema debe estandarizar pa\'ises, variables, unidades,
monedas, fechas, frecuencias y horizontes mediante un modelo com\'un.}

\rf{norm-formatos}{El sistema debe transformar informaci\'on publicada en API,
Excel, CSV, PDF, HTML y JSON hacia el modelo com\'un.}

\rf{norm-historico}{El sistema debe reconstruir series hist\'oricas, conservar
las revisiones sucesivas de cada dato y mantener el estado de la informaci\'on
disponible en cada momento.}

\subsubsection{Control de calidad y trazabilidad}

\rf{qa-deteccion}{El sistema debe detectar duplicados, faltantes,
inconsistencias y cambios de formato antes de publicar informaci\'on.}

\rf{qa-cuarentena}{El sistema debe enviar los registros dudosos a cuarentena,
conservar los originales inmutables en la capa Bronze y registrar ejecuciones,
errores y correcciones.}

% ---------------------------------------------------------------------
\subsection{Capa de inteligencia}

Este bloque cubre lo que el sistema calcula sobre los datos ya normalizados: las
métricas propias, las comparaciones entre países y las señales analíticas
derivadas de ambas.

\subsubsection{Métricas}

\rf{met-calculo}{El sistema debe calcular cambios, momentum, divergencias
regionales, percentiles hist\'oricos, sorpresas y persistencia sobre las series
normalizadas.}

\rf{met-config}{El sistema debe permitir configurar ventanas, par\'ametros y
metodolog\'ias de c\'alculo seg\'un cada m\'odulo.}

\rf{met-metas-cruces}{El sistema debe calcular la distancia frente a metas y
resolver los cruces entre expectativas, inflaci\'on observada y actividad
efectiva.}

\subsubsection{Comparación regional}

\rf{comp-paises}{El sistema debe permitir comparar pa\'ises, variables,
per\'iodos y m\'odulos econ\'omicos de Am\'erica Latina bajo criterios
homog\'eneos.}

\rf{comp-rankings}{El sistema debe permitir identificar revisiones, divergencias
y \emph{rankings} regionales relevantes.}

\rf{comp-advertencia}{El sistema debe advertir al usuario cuando un cruce no
pueda calcularse, distinguiendo si la causa es la falta de datos o la ausencia
de los m\'odulos requeridos en el plan contratado.}

\subsubsection{Señales analíticas}

\rf{senal-alertas}{El sistema debe generar alertas e informes descriptivos sobre
cambios relevantes previamente calculados y validados.}

\rf{senal-clasificacion}{El sistema debe clasificar los movimientos seg\'un su
direcci\'on, intensidad y relevancia para el an\'alisis.}

\rf{senal-resumen}{El sistema debe resumir qu\'e se movi\'o en la regi\'on sin
que el usuario deba descargar, limpiar o reconstruir manualmente los datos. El
resumen debe derivarse exclusivamente de m\'etricas ya calculadas y validadas.}

% ---------------------------------------------------------------------
\subsection{Aplicación web}

Este bloque cubre aquello con lo que el usuario interactúa directamente: el
acceso a su cuenta, el control de permisos según lo contratado, la consulta y
visualización de la información, y las vías por las que se lleva el resultado
fuera de la plataforma.

\subsubsection{Autenticación y usuarios}

\rf{auth-registro}{El sistema debe permitir a los usuarios registrarse e iniciar
sesi\'on con sus credenciales.}

\rf{auth-recuperacion}{El sistema debe permitir al usuario recuperar el acceso a
su cuenta.}

\rf{auth-entitlements}{El sistema debe aplicar roles, permisos y
\emph{entitlements} seg\'un el plan, los m\'odulos, los pa\'ises y las
funcionalidades contratadas.}

\rf{auth-admin}{El sistema debe permitir al administrador gestionar usuarios,
organizaciones, planes, m\'odulos, permisos y suscripciones.}

\subsubsection{Pagos y suscripciones}

\rf{pago-checkout}{El sistema debe delegar el pago al \emph{checkout} de un
proveedor externo. \proyectoNombre{} no almacena ni procesa directamente datos
de tarjetas.}

\rf{pago-webhook}{El sistema debe recibir y procesar las notificaciones del
proveedor de pagos cuando una suscripci\'on sea aprobada, renovada, rechazada,
cancelada o reembolsada.}

\rf{pago-datos}{El sistema debe conservar \'unicamente los identificadores
externos, el estado de la suscripci\'on y los datos necesarios para actualizar
m\'odulos, planes y permisos.}

\rf{pago-cambio-plan}{El sistema debe reflejar los cambios de plan en los
\emph{entitlements} de la cuenta, habilitando o revocando el acceso a m\'odulos y
funcionalidades seg\'un corresponda.}

\subsubsection{Dashboard de análisis}

\rf{dash-seleccion}{El sistema debe permitir al usuario seleccionar pa\'is,
m\'odulo, variable, fuente y per\'iodo hist\'orico.}

\rf{dash-visualizacion}{El sistema debe permitir visualizar gr\'aficos, tablas,
m\'etricas, \emph{rankings} y cruces entre m\'odulos desde una misma interfaz.}

\rf{dash-navegacion}{El sistema debe permitir navegar entre pa\'ises, variables,
m\'odulos y a\~nos de forma simple y r\'apida.}

\rf{dash-historico}{El sistema debe permitir consultar hist\'oricos y revisiones
para analizar cambios y tendencias en el tiempo.}

\rf{dash-procedencia}{El sistema debe mostrar, junto a cada dato y m\'etrica
visualizada, su fuente, fecha, versi\'on y nota metodol\'ogica.}

\subsubsection{Radar de expectativas}

\rf{radar-acceso}{El sistema debe publicar un radar peri\'odico con los
principales cambios econ\'omicos de Am\'erica Latina.}

\rf{radar-free}{El sistema debe mostrar una versi\'on resumida del radar a los
usuarios sin suscripci\'on vigente, sin requerir contrataci\'on.}

\rf{radar-pago}{El sistema debe mostrar la versi\'on completa del radar, las
alertas y los \emph{briefings} a los usuarios cuya suscripci\'on habilite esas
funcionalidades.}

\subsubsection{Exportación y uso de resultados}

\rf{exp-resultados}{El sistema debe permitir exportar resultados transformados
para an\'alisis, estudio o reportes a los usuarios cuya suscripci\'on habilite
la exportaci\'on.}

\rf{exp-formatos}{El sistema debe permitir descargar informaci\'on en formatos
de uso corriente, como Excel y PDF.}

\rf{exp-trazabilidad}{El sistema debe incluir en cada exportaci\'on la fuente,
fecha, versi\'on y nota metodol\'ogica de los datos y m\'etricas exportados.}

\rf{exp-consultas}{El sistema debe permitir guardar consultas frecuentes para
facilitar el seguimiento recurrente de expectativas.}

\subsubsection{Integración programática y fuentes privadas}

\rf{api-consulta}{El sistema debe permitir a las instituciones consultar datos y
m\'etricas mediante una API autenticada, con cuotas, l\'imites y permisos seg\'un
el plan contratado.}

\rf{api-privadas}{El sistema debe permitir a las instituciones integrar fuentes
privadas dentro de un entorno aislado, sin mezclarlas con la informaci\'on de
otros clientes.}

% =====================================================================
\section{Requerimientos no funcionales}
\label{sec:rnf}

Los requerimientos no funcionales describen cómo debe comportarse el sistema al
cumplir lo anterior: con qué nivel de trazabilidad, seguridad, aislamiento y
cumplimiento normativo. No agregan funcionalidad visible, pero condicionan las
decisiones de arquitectura: cada decisión documentada en las Notas de
Arquitectura y Diseño debe indicar a cuál de estos responde.

Los umbrales medibles de calidad ---frescura del dato, autonomía de la ingesta,
cobertura de trazabilidad, tasa de cuarentena, tiempo de respuesta y
disponibilidad--- se incorporarán a esta sección como requerimientos numerados
una vez que queden definidos en el Documento de Visión.

\subsection{Trazabilidad y auditoría}

\rnf{traza-completa}{Cada dato, m\'etrica y exportaci\'on debe conservar su
fuente, fecha, versi\'on y criterio de procesamiento.}

\rnf{traza-inmutable}{Los archivos y respuestas originales de las fuentes deben
conservarse inmutables en la capa Bronze.}

\rnf{traza-auditoria}{El sistema debe mantener registros de auditor\'ia.}

\subsection{Seguridad}

\rnf{seg-superficie}{La plataforma debe exponer \'unicamente el frontend y los
\emph{endpoints} necesarios.}

\rnf{seg-sesiones}{El sistema debe utilizar sesiones o \emph{tokens} seguros y
validar las solicitudes recibidas.}

\rnf{seg-secretos}{La gesti\'on de secretos debe realizarse fuera del c\'odigo
fuente.}

\rnf{seg-permisos}{El sistema debe aplicar roles y permisos por usuario,
organizaci\'on, plan, m\'odulo y pa\'is.}

\rnf{seg-pci}{\proyectoNombre{} no debe almacenar ni procesar directamente datos
de tarjetas.}

\subsection{Aislamiento}

\rnf{multi-aislamiento}{El sistema debe proveer aislamiento
multi-\emph{tenant} para las fuentes privadas.}

\subsection{Cumplimiento}

\rnf{cump-transformado}{El producto debe entregar informaci\'on transformada
---m\'etricas, comparaciones y valor agregado--- y no una copia del dato crudo.}

\rnf{cump-atribucion}{Cada visualizaci\'on, exportaci\'on o respuesta de API
debe incluir fuente, fecha, versi\'on y nota metodol\'ogica.}

\end{document}
